\documentclass[11pt,a4paper]{article}
\usepackage[margin=25mm]{geometry}
\usepackage{amsmath,amssymb,amsthm,mathtools}
\usepackage{booktabs,longtable,array,enumitem,microtype}
\usepackage{xcolor}
\usepackage[colorlinks=true,linkcolor=blue!45!black,citecolor=blue!45!black,urlcolor=blue!45!black]{hyperref}
\usepackage{fancyhdr}
\pdfmapfile{+cm.map}
\pdfmapfile{+cmextra.map}
\pdfmapfile{+symbols.map}
\pdfmapfile{+latxfont.map}
\setlength{\emergencystretch}{3em}
\setlength{\parskip}{4pt}
\setlength{\parindent}{0pt}
\setlist{itemsep=3pt,topsep=5pt,leftmargin=*}
\pagestyle{fancy}\fancyhf{}\fancyhead[L]{\small GCT\_NSRH: checked reductions and open gates}\fancyfoot[C]{\thepage}
\setlength{\headheight}{14pt}
\newcommand{\Q}{\mathbb Q}\newcommand{\Z}{\mathbb Z}\newcommand{\C}{\mathbb C}\newcommand{\N}{\mathbb N}\newcommand{\F}{\mathbb F}
\newcommand{\A}{\mathbb A}\newcommand{\PP}{\mathbb P}\newcommand{\GL}{\mathrm{GL}}
\newcommand{\Hom}{\operatorname{Hom}}\newcommand{\Spec}{\operatorname{Spec}}\newcommand{\Tr}{\operatorname{Tr}}\newcommand{\rank}{\operatorname{rank}}\newcommand{\Sym}{\operatorname{Sym}}\newcommand{\Res}{\operatorname{Res}}\newcommand{\IC}{\operatorname{IC}}\newcommand{\Gr}{\operatorname{Gr}}\newcommand{\Rep}{\operatorname{Rep}}\newcommand{\Frob}{\operatorname{Frob}}\newcommand{\End}{\operatorname{End}}
\newcommand{\id}{\operatorname{id}}\newcommand{\et}{\mathrm{\acute et}}\newcommand{\boxt}{\boxtimes}
\newcommand{\fid}[1]{\textsf{#1}}\newcommand{\gate}[1]{\textsf{#1}}
\newtheoremstyle{named}{8pt}{5pt}{\itshape}{} {\bfseries}{.}{.5em}{}
\theoremstyle{named}\newtheorem{finding}{Finding}[section]
\theoremstyle{definition}\newtheorem{definition}[finding]{Definition}
\theoremstyle{remark}\newtheorem{remark}[finding]{Remark}
\title{Purity from path sums and the nonstandard GCT positivity problem\\[4pt]\large Checked reductions, explicit barriers, and the missing realization}
\author{Research run \texttt{GCT\_NSRH\_20261121}\thanks{Prepared by GPT-6 Astra Pro with three parallel source/proof reviews. Model identity is supplied by the runtime; the precise reasoning-effort setting is not exposed. Check types and actual scopes are recorded in \texttt{CHECK\_EVENTS.tsv}. No proof-assistant verification of these findings is claimed.}}
\date{21 November 2026 (Europe/Berlin)}
\begin{document}\maketitle
\begin{abstract}
We investigate whether the path-quadric sign mechanism described in Alp\"oge's supplied \emph{Remark} transfers to Mulmuley's nonstandard quantum groups. We prove an exact vector-valued path-sum estimate and a conditional purity criterion. For trace functions on $\A^d/\F_q$ bounded in pointwise weight $a$, iterated Parseval gives an error of size $q^{n(k(a+d)/2+d/2+1)}$. If the relevant spectra have opposite cohomological parity for unbounded path lengths $k$, positive spectral multiplicities prevent cancellation. Taking first extension degree $n$ and then $k$ to infinity gives the upper weight $a+d$; a specified Tate duality gives equality. An explicit sublinear boundary-error version is included.

The supplied GCT papers do not formulate a finite-field nonstandard Riemann hypothesis with designated Frobenius objects and weights. Their positivity targets also include signed factors, rational normalizations, and, in general plethysm, corrections to positivity. We isolate the additional realization, effective-category, basis, Lefschetz, specialization, and complexity hypotheses. A rank-one calculation in the actual nonstandard Hecke inclusion proves that its ambient Kazhdan--Lusztig positive cone meets the nonstandard algebra only in scalar multiples of the unit. Thus a straightforward inheritance of the standard effective category fails even in this small case. Pure two-point geometry supplies a separate counterexample to determinant positivity implying positive structure constants in a chosen effective basis.

Theta/interlacing arguments yield an exact conditional path-count theorem, but no functor identifying those paths with the required GCT multiplicity spaces is supplied. Existing rectangular positivity and no-occurrence theorems are stated with their parameters. All checked findings are translated into schemes, representation $K_0$, $\lambda$-operations and Weil realizations, with verified archival anchors.
\end{abstract}
\begin{center}\fbox{\parbox{.91\linewidth}{\textbf{Final research status: INCOMPLETE.} The conditional theorems and counterexamples below are proved and checked. No nonstandard RH for the designated GCT objects, general canonical-basis positivity, general positive plethysm rule, or complexity separation is proved. The final part records the exact open gates.}}
\end{center}
\tableofcontents\newpage

\part{The exact question and its sources}
\section{Scope, premises, and conventions}\label{sec:scope}
The main question is a transfer question. We accept the RH and Dirichlet-GRH results in the supplied Alp\"oge material as stipulated premises, as requested. We do not re-referee those results. Their acceptance does not supply a theorem relating an automorphic or finite-field trace system to a nonstandard quantum canonical basis.

The relevant finite-field description is in \cite[printed p.~3, footnote~3, attached to the paragraph before the timeline; also step~3]{Alpoge}. Timeline step~5 itself concerns a von Mangoldt functional and continues onto p.~4. We use ``the step-5 mechanism'' in the prompt's sense, but retain this more precise locator. The note names a path-quadric identity, middle-cohomological signs, iterated Parseval, and the absence of cancellation of inverse roots. It does not state a standalone theorem with all the geometric hypotheses used below.

The supplied \texttt{ALPOGE\_Main.lean} consists of five exported wrappers, including \texttt{riemannHypothesis := rhOfAll}, and imports \texttt{Supplements.Kernels.MainTheoremsAll}. The supplied \texttt{ALPOGE\_Statements.lean} defines the zeta/Dirichlet zero and nonvanishing propositions. Neither file contains a nonstandard quantum-group, path-concentration, or theta/GCT transfer theorem. The imported development is not in this ZIP. The supplied scan records zero occurrences for the listed prohibited constructs; this run did not repeat that scan or execute the missing development. These are access and locator facts, not objections to the stipulated theorem. No Masullo result is used: none is identified or supplied in this prompt.

We write $v$ for the \emph{formal quantum parameter}, $q$ for a fixed prime power, and $Q=q^n$ for an extension-field cardinality. In the GCT sources the quantum parameter is usually called $q$. No substitution $v=Q$ or $v=Q^{1/2}$ constructs a Weil structure. All weights in the purity conclusion are integers. The Frobenius convention is geometric Frobenius, so $\Q_\ell(-1)$ has eigenvalue $q$ over $\F_q$.

``Finding'' means a deduction established for this investigation, including elementary lemmas in the precise form needed. It is not a claim of priority over the mathematical literature. Known research results are cited, not presented as new findings. A source statement checked in context is not thereby an independently checked proof of the entire cited theorem.

\section{What IV, VI, VII and VIII actually specify}\label{sec:sources}
\subsection{Versions and the missing RH statement}
The attached IV is the revised Blasiak--Mulmuley--Sohoni paper, arXiv cs/0703110v4 (2013), not merely the 2007 version referred to by VII and VIII. It proves nonstandard Schur--Weyl duality and the two-row crystal result. VI is arXiv 0704.0229v4 (2009); VII and VIII are 0709.0749v2 and 0709.0751v2 (2008). Rendered title-page dates produced by an undated \verb|\maketitle| are not used as version dates.

\textbf{Source finding A01.} In these four supplied versions there is no formulated nonstandard RH specifying finite-field objects, Frobenius actions, cohomological degrees and weights. This is an exact-source conclusion, not a theorem about what other work might contain. The relevant locators are:

{\small\begin{longtable}{p{.13\linewidth}p{.35\linewidth}p{.44\linewidth}}
\toprule Work&Exact locator&Content and status\\\midrule\endhead
IV&\S1.1, Problem 1.1; \S\S1.2--1.5; \S19&Positive-rule and basis program; the revised text has no finite-field purity statement.\\
VI&Introduction, Figure 1.1; \S1.7, Observation 1.7.1 and Figure 1.2&The proposed RH/extensions connection is marked by question marks; standard canonical positivity motivates the program.\\
VI&\S4.2, Proposition 4.2.1; \S4.2.2, Question 4.2.6 (PH0??); \S4.3&Asks for compatible canonical bases and analogous positivity. No nonstandard weight theorem is stated.\\
VII&\S1 after Conjectures 1.3--1.4; \S4.1, equations (9)--(10), Theorem 4.4; \S4.2, Theorem 4.5&Nonstandard minors and Laplace coefficients need interpretations analogous to perverse-sheaf constructions.\\
VIII&\S1, pp.~2--3; \S2.3 after Conjecture 2.5&An extension of perverse sheaves and finite-field RH is proposed as a way to understand minors and globalization.\\
VIII&\S2.4.2 after Conjecture 2.7; \S2.4.3 after Conjecture 2.9&Topological interpretations, Hard Lefschetz, and a theory for noncommutative, nonflat objects are suggested; the general case includes correction terms.\\\bottomrule
\end{longtable}}
These statements are supported by complete reading of their contexts and matching numbered PDF statements, not merely a search for the phrase ``Riemann hypothesis.'' Their appropriate conclusion is that a precise realization/purity package must first be specified.

\subsection{The algebraic objects}
Let $H$ be a complex connected classical reductive group, $X=V_\mu(H)$, $G=\GL(X)$, and $\rho:H\to G$ the representation homomorphism. The generalized plethysm multiplicity is
\begin{equation}\label{eq:plethysm}
a_{\lambda,\mu}^{\pi}=\dim\Hom_H\bigl(V_\pi,\Res_\rho V_\lambda(G)\bigr).
\end{equation}
For Kronecker coefficients, $H=\GL(V)\times\GL(W)$, $X=V\otimes W$, and $\rho(g,h)=g\otimes h$.

VII, Problem~1.1 and Theorems~1.2, 2.1 and 3.1, constructs a nonstandard matrix-coordinate bialgebra
\[
 S=\mathcal O(M_v^H(X)),\qquad \mathcal A_r^H=(S_r)^*,
\]
with a quantized map from $H_v$ to the associated nonstandard quantum object. Its matrix generators satisfy the projector relation
\[
 P^{+,H}_{X,X}(\mathbf u\otimes\mathbf u)
 =(\mathbf u\otimes\mathbf u)P^{+,H}_{X,X}
\]
(VII, \S2, equations (4)--(6)). The projector is constructed from the braiding on $X_v\otimes X_v$. The determinant may vanish, so the strict term may be quantum semigroup; an antipode is not automatic. The dual nonstandard symmetric-group deformation is $\mathcal B_r^H(v)$ (VII, \S5).

These are noncommutative algebras over fields containing $\Q(v)$, sometimes an algebraic extension. Their bar involution is not arithmetic Frobenius. Complete reducibility or a compact real form is not a Weil structure. Nonflatness and noncommutativity have different roles: nonflat commutative schemes still have \(\ell\)-adic sheaves, whereas a noncommutative coordinate algebra is not the coordinate ring of a scheme in the EGA/SGA sense.

General plethysm duality and reciprocity are VII, Conjectures~1.3--1.4 and refined Conjecture~6.1. They would give
\begin{equation}\label{eq:comparison}
a_{\lambda,\mu}^{\pi}=\sum_\alpha m_\lambda^\alpha n_\pi^\alpha,
\end{equation}
where $n_\pi^\alpha$ is an $H_v$ constituent multiplicity and $m_\lambda^\alpha$ is a Specht multiplicity after specialization at $v=1$. In revised IV, Theorems~1.5/12.1 establish the Kronecker duality, and \S1.4, equation~(11), gives the corresponding Kronecker decomposition. It would be incorrect to leave that particular comparison conjectural merely because VII predates the revision.

\subsection{The actual positivity targets}\label{sec:targets}
The basis constructions themselves have separate prerequisites. VIII, Conjecture~2.2, concerns compatible local crystal bases; Question~2.3's ordinary balanced-triple condition fails in examples; Conjecture~2.5 proposes a unique minimum-degree globalization. Conjectures~2.6 and 3.3 concern cells, Conjectures~3.1--3.2 the dual globalization, and Conjecture~3.7 the \emph{opposite} basis's quasi-subcells and specialization. VIII's Proposition~2.4 contains an omitted proof and assumes the local-crystal conjecture; we do not record it as independently proved.

\paragraph{Kronecker, coordinate side.}
VIII, Conjecture~2.7, asks that expansion coefficients of $CB(\mathcal A_r^H)$ in the external canonical basis $CB(Z)$, $Z=\End(X_v^{\otimes r})$, be polynomials with nonnegative \emph{rational} coefficients (equation~(9)). Multiplication and comultiplication constants have the form
\begin{equation}\label{eq:kroneckerpositive}
 \pm(v-v^{-1})^a f(v),\qquad a\in\Z_{\ge0},
\end{equation}
where $f$ is bar-invariant, coefficientwise nonnegative rational and unimodal. The same is proposed for $CB(S)$.

\paragraph{Kronecker, dual side.}
VIII, Conjecture~3.4, applies to structure constants and external expansion coefficients of $E(r)$ and $E^{\mathrm{opp}}(r)$. Its form is
\begin{equation}\label{eq:dualpositive}
 \pm(v^{1/2}-v^{-1/2})^a f(v),
\end{equation}
with $f$ nonnegative rational, bar-invariant and unimodal in the half-power variable. Thus raw coefficientwise positivity is not the stated assertion even in the Kronecker case.

\paragraph{General plethysm.}
VIII, Conjectures~2.8 and 3.5, concern normalized coefficients of minimal polynomials of algebraic structure/expansion coefficients. The displayed normalizations are equations~(10) and (14):
\[
 (-1)^e(v-v^{-1})^{e'}c(v),\qquad
 (-1)^e(v^{1/2}-v^{-1/2})^{e'}g(v).
\]
The nonzero normalized polynomial is required to be positive at $v=1$ (called saturated there). Conjectures~2.9 and 3.6 ask for a dominant positive unimodal part and a small correction, for example
\[
 c=c^0+c^1,\qquad
 \frac{\|c^1\|_2}{\|c^0\|_2}
 \le\frac1{\operatorname{poly}(\langle\mu\rangle,\langle\pi\rangle,\langle r\rangle)},
 \qquad c^0(1)\gg |c^1(1)|.
\]
Conjecture~3.6 is explicitly informal. A theorem target must fix the quantitative meaning of the polynomial and the asymptotic comparison. VIII, \S5.4.1, Figure~30, displays negative entries in a normalized example; these are cited source computations, not a recomputation of the whole algebra in this run.

There is also a sign-convention issue in ``coefficients of the minimal polynomial'': the usual monic polynomial of 1 is $T-1$, while the source later refers to trace and norm without the conventional sign on trace. We record this ambiguity rather than use it to disprove the intended conjecture. ``Frobenius automorphisms'' in VIII, \S5.4.1, refers to algebraic conjugations over $\Q(v)$, not to finite-field Weil Frobenius.

\subsection{Positivity hypotheses in the flip}
VI, Hypothesis~1.6.4 (PH1), is substantially stronger than numerical nonnegativity. For each triple there should be a rational polytope $P_{\lambda,\mu}^{\pi}$ whose Ehrhart function is
\[
 \widetilde a_{\lambda,\mu}^{\pi}(t)=a_{t\lambda,\mu}^{t\pi},
\]
with $\mu$ fixed (VI, \S1.6, definitions preceding Theorem~1.6.1). It must admit a description $Ax\le b$ with $A$ independent of $\lambda,\pi$ and $b$ homogeneous linear in them; ambient dimension and entry bitlengths are polynomially bounded as specified there; membership and a separating hyperplane are computable in polynomial time. The number of inequalities need not be polynomial.

VI, Hypothesis~1.6.5 (SH), asks for a polynomial bound on the saturation index in $\dim G+\operatorname{ht}\lambda+\operatorname{ht}\pi$ and asymptotic generic strict saturation. Hypothesis~1.6.6 (PH2) bounds the positivity index; Hypothesis~1.6.7 (PH3) bounds the modular index of a positive rational Hilbert-series form. Strict saturation in every case is not the attached revised hypothesis.

The conditional implication to relaxed PHflip is VI, Theorem~1.6.8. A printed discrepancy must be retained: its part~(a) displays $a_{b\lambda,b\mu}^{b\pi}$ although the stretching definition and its stated proof through PH1 keep $\mu$ fixed. We use only the fixed-$\mu$ consequence directly justified by the displayed Ehrhart model and record the all-three-scaling formulation as requiring clarification. This editorial point is not needed for our main obstruction. More fundamentally, VIII, \S6, explicitly asks for additional effective-label properties. An abstract canonical basis or an abstract positive count does not prove PH1, SH, or the existence of a lower-bound obstruction.

\section{The Kazhdan--Lusztig template, with weights fixed}\label{sec:KL}
Let $G/\F_Q$ be split reductive, $B$ a Borel subgroup, $C_w=BwB/B$ and $X_w=\overline{C_w}$. Put
\[
 \IC_w^u=j_{!*}\Q_\ell[\ell(w)][-\ell(w)]
\]
with the usual extension to $X_w$; the notation records the unshifted convention, constant in degree zero on $C_w$. Then
\begin{equation}\label{eq:KL}
 P_{x,w}(T)=\sum_{j\ge0}\dim\mathcal H^{2j}(\IC_w^u)_x\,T^j.
\end{equation}
Odd stalk groups vanish, and each eigenvalue in degree $2j$ is $Q^j$. This is the Tate/parity statement, stronger than an absolute-value bound. Exact primary locators for this form are Haines \cite[Theorems 2.0.1, 5.0.4; Proposition 3.0.2; \S4]{Haines}. The original KL geometric identification is \cite[Theorem 4.3]{KL}.

In the modern proof, a Bott--Samelson resolution is proper, its fibers admit affine pavings, and the decomposition theorem isolates the intersection complex. The affine paving provides even/Tate cohomology. The compatibility with Weil structures must be checked: a geometric splitting need not itself be Frobenius-equivariant. Haines \S4 discusses precisely this issue. His introduction also notes that the original KL proof does not require the decomposition theorem; the latter is the requested modern template, not a claim about historical logical necessity.

Deligne's upper-weight theorem is \cite[Theorem 3.3.1; Corollaries 3.3.4--3.3.6]{WeilII}. BBD intermediate extension preserves purity (5.3.2), pure perverse sheaves are geometrically semisimple (5.3.8), and pure complexes split geometrically (5.4.5); relative Hard Lefschetz is 5.4.10 \cite{BBD}. None of these statements identifies an arbitrary pure trace with the dimension polynomial \eqref{eq:KL}. Arithmetic Frobenius semisimplicity is distinct: BBD, pp.~138--139, explicitly discusses pure unipotent Jordan-block sheaves.

{\small\begin{longtable}{p{.25\linewidth}p{.29\linewidth}p{.38\linewidth}}
\toprule Standard ingredient&Nonstandard datum supplied&Unfilled role\\\midrule\endhead
Flag schemes, Schubert strata&$H_v$, nonstandard FRT bialgebra; support-defined cells&An actual space/category with those strata and stalk functors.\\
Distinguished $\IC_w$ and its normalized class&Candidate minimum-degree bases&Objects whose classes are the \emph{specified} bases; existence and uniqueness.\\
Weil structure and closed-point Frobenius&Algebraic bar involution&Arithmetic models, traces, degrees and integer weights.\\
Proper resolutions/convolution&Algebra multiplication and comultiplication&Correspondences or functors with a decomposition theorem.\\
Tate parity and effective multiplicities&Signed rational forms in VIII&Positive multiplicity objects, signs, factors and denominators.\\
Verdier duality and Hard Lefschetz&Bar symmetry/unimodality conjectures&Compatible duality and Lefschetz isomorphisms.\\
Crystal/path labels&Conjectural cells; proved two-row tableaux&Effective labels, specialization and polynomial-time verification in general.\\\bottomrule
\end{longtable}}
For standard quantum groups, quiver and flag constructions supply related geometric realizations; it would be misleading to assert that every quantum canonical-basis constant is a stalk of one Schubert variety.

\part{Checked findings}
\section{The sign mechanism and its exact strength}\label{sec:main}
This section contains the main proved reduction. Its hypotheses explicitly retain the information that the phrase ``middle cohomology'' compresses.

\begin{finding}[\fid{F01}: path identity and Parseval estimate]\label{f:path}
Let $V=\F_Q^d$, $d\ge1$, with a nondegenerate bilinear pairing, and let $\psi$ be a nontrivial additive character. For $k\ge2$ and any function $t:V\to\C$, put
\[
 M=\sum_{x\in V}t(x),\quad
 P_k(x_1,\ldots,x_k)=\sum_{i=1}^{k-1}\langle x_i,x_{i+1}\rangle,
 \quad N_k=\sum_{P_k=0}\prod_{i=1}^k t(x_i).
\]
Then $QN_k=M^k+R_k$, where
\[
 R_k=\sum_{h\in\F_Q^\times}\sum_{x_1,\ldots,x_k\in V}
       \prod_i t(x_i)\,\psi(hP_k),
\]
and
\begin{equation}\label{eq:parseval}
 |R_k|\le (Q-1)Q^{d(k-1)/2}\|t\|_2^2\|t\|_\infty^{k-2}.
\end{equation}
No nonnegativity assumption on $t$ is required.
\end{finding}
\begin{proof}
Orthogonality gives $\sum_h\psi(hz)=Q\,1_{z=0}$; separating $h=0$ gives the identity. For $h\ne0$, set $K_h(x,y)=\psi(h\langle x,y\rangle)$ and $D=\operatorname{diag}(t)$. Nondegeneracy gives
\[
 (K_hK_h^*)_{x,x'}=\sum_y\psi(h\langle x-x',y\rangle)
 =Q^d1_{x=x'},
\]
so $\|K_h\|=Q^{d/2}$. The inner sum is
$t^{\mathsf T}K_hD K_h\cdots D K_ht$, with $k-1$ kernels and $k-2$ diagonal factors. The norm of the row functional $t^{\mathsf T}$ is $\|t\|_2$, even for complex $t$ (it is the Hermitian dual of $\bar t$). Cauchy--Schwarz and submultiplicativity give the bound for one $h$; sum the $Q-1$ terms. For $k=2$ there are no diagonal factors.
\end{proof}
This refines Fable suggestion~3's scalar toy in the exact form needed here. It is an open-path matrix coefficient, not in general a trace moment: tracing the power closes the path into a cycle.

\begin{finding}[\fid{F02}: finite spectral no-cancellation]\label{f:spectral}
For distinct nonzero complex numbers $\alpha_j$ and nonzero constants $c_j$, let $p_n=\sum_jc_j\alpha_j^n$. For a nonempty finite spectrum,
\[
 \limsup_{n\to\infty}|p_n|^{1/n}=\max_j|\alpha_j|.
\]
Consequently $|p_n|\le Cn^A R^n$ for all $n\ge1$, $R>0$, implies $|\alpha_j|\le R$. The same holds for any subexponential prefactor.
\end{finding}
\begin{proof}
Near zero,
\[
 \sum_{n\ge1}p_nz^n=\sum_j\frac{c_j\alpha_jz}{1-\alpha_jz}.
\]
Every $1/\alpha_j$ is a genuine pole, since the other summands are regular there. The radius of convergence is $1/\max|\alpha_j|$. Cauchy--Hadamard proves the equality; the bound gives radius at least $1/R$. The empty case is vacuous. Coincident spectral values must first be grouped; the conclusion concerns their \emph{nonzero} total coefficients.
\end{proof}
Positive multiplicities are one sufficient way to ensure that regrouping does not erase a root. Eigenvalues themselves need not be positive, or real. The traces of powers do not see nilpotent Jordan off-diagonal entries, so semisimplicity is not required for this lemma.

\begin{finding}[\fid{F03}: tensor-path purity criterion]\label{f:criterion}
Fix a prime power $q>1$, an integer $a$, $d\ge1$, and $B>0$. For every $n\ge1$ let $Q=q^n$ and $t_n:\F_Q^d\to\C$ satisfy $|t_n(x)|\le BQ^{a/2}$. Form $M_n,N_{k,n},R_{k,n}$ as in \fid{F01}. Assume:
\begin{enumerate}[label=(S\arabic*)]
\item A fixed finite multiset $\mathcal A$ of nonzero algebraic numbers, with positive integer multiplicities, and $\epsilon\in\{1,-1\}$ satisfy
\[
 M_n=\epsilon\sum_{\alpha\in\mathcal A}\alpha^n\qquad(n\ge1).
\]
\item For an unbounded set of integers $k\ge2$ there is, for each $k$, a fixed finite multiset $\mathcal B_k$ of nonzero algebraic numbers, possibly empty, with positive integer multiplicities such that
\[
 q^nN_{k,n}=-\epsilon^k\sum_{\beta\in\mathcal B_k}\beta^n\qquad(n\ge1).
\]
\end{enumerate}
Then $|\alpha|\le q^{(a+d)/2}$ for every $\alpha\in\mathcal A$.

For the all-conjugate assertion, require for every embedding $\sigma$ of the algebraic spectral data a corresponding trace system $t_{n,\sigma}$ satisfying these hypotheses with spectra $\sigma(\mathcal A),\sigma(\mathcal B_k)$ and a constant $B_\sigma$ independent of $n$. If in addition
\begin{equation}\label{eq:duality}
 \mathcal A=q^{a+d}/\mathcal A
 \quad\text{as multisets},
\end{equation}
then every eigenvalue has every complex absolute value $q^{(a+d)/2}$: the spectrum is pure of weight $a+d$.
\end{finding}
\begin{proof}
The pointwise bound gives $\|t_n\|_2^2\le B^2Q^{a+d}$. Equation~\eqref{eq:parseval} therefore gives
\begin{equation}\label{eq:overhead}
 |R_{k,n}|\le B^kQ^{,k(a+d)/2+d/2+1}.
\end{equation}
By the exact path identity and (S1)--(S2),
\begin{equation}\label{eq:positivespectrum}
 -\epsilon^kR_{k,n}=
 \sum_{\beta\in\mathcal B_k}\beta^n+
 \sum_{(\alpha_1,\ldots,\alpha_k)\in\mathcal A^k}
 (\alpha_1\cdots\alpha_k)^n.
\end{equation}
After grouping equal roots all coefficients are positive integers. In particular the diagonal tuple contributes $\alpha^k$ and cannot be cancelled. For each \emph{fixed} $k$, \fid{F02} and \eqref{eq:overhead} give
\[
 |\alpha|^k\le q^{k(a+d)/2+d/2+1},\qquad
 |\alpha|\le q^{(a+d)/2+(d/2+1)/k}.
\]
Let $k$ tend to infinity through the unbounded set. This order of limits is essential: constants such as $B^k$ disappear in the $n$th-root limit before $k$ varies. Applying the same argument to each $\sigma$ gives all-conjugate upper bounds. Since $a+d$ is an integer, $q^{a+d}$ is rational and fixed by every $\sigma$. Apply the upper bound to $q^{a+d}/\sigma(\alpha)$, which occurs by \eqref{eq:duality}, to obtain the reverse inequality.
\end{proof}
The purity conclusion is conditional on a \emph{specified Tate duality}. An early draft allowed an arbitrary algebraic dualizing scalar; the independent check found that its conjugates could have different absolute values. The final theorem uses the integer-weight scalar $q^{a+d}$ and the correction is recorded in the audit.

\begin{finding}[\fid{F04}: controlled background terms]\label{f:errors}
Replace (S2) by
\[
 q^nN_{k,n}=-\epsilon^k\sum_{\beta\in\mathcal B_k}\beta^n+E_{k,n}.
\]
Suppose $|E_{k,n}|\le C_kn^{A_k}q^{n(k(a+d)/2+\eta_k)}$ for all $n$, with $\eta_k/k\to0$ along the unbounded set of $k$. Then \fid{F03}'s upper bound, and its purity conclusion under the same all-embedding duality hypotheses, still hold.
\end{finding}
\begin{proof}
The positive spectral sum in \eqref{eq:positivespectrum} is now $\epsilon^k(E_{k,n}-R_{k,n})$. Its exponential rate for fixed $k$ is at most
$q^{k(a+d)/2+\max(d/2+1,\eta_k)}$. Apply \fid{F02} and take the $k$th root. Since $\max(d/2+1,\eta_k)/k\to0$, the same limit follows. Duality supplies the reverse bound as before.
\end{proof}
This is a sufficient error budget, not a necessity theorem for every possible proof. It permits suitable removed Tate or boundary terms; it does not justify discarding uncontrolled terms.

\subsection{What the hypotheses mean geometrically}\label{sec:tracegeom}
Let $\mathcal F$ be a constructible Weil sheaf on $\A^d/\F_q$. Its trace functions $t_n$ have the required pointwise estimate if stalk ranks are bounded by $B$ and all stalk eigenvalues have weight at most $a$. Let
\[
 H_k=\{(x_1,\ldots,x_k):P_k(x_1,\ldots,x_k)=0\}
 \subset(\A^d)^k.
\]
The character kernel is represented by an Artin--Schreier sheaf pulled back along $hP_k$; the coefficient system on $H_k$ is the restriction of $\mathcal F^{\boxt k}$. The Grothendieck trace formula identifies the sums with alternating traces of $R\Gamma_c$.

If $R\Gamma_c(\A^d_{\bar{\F}_q},\mathcal F)$ is concentrated in degree $d$, then $\epsilon=(-1)^d$. If the relevant cohomology of $H_k$ is concentrated in degree $kd-1$, its sign is
\[
 (-1)^{kd-1}=-\bigl((-1)^d\bigr)^k.
\]
Multiplication by $q^n$ is the twist $(-1)$. These facts give precisely (S1)--(S2), or \fid{F04} when background terms have been controlled. A suitable duality of the coefficient system and cohomology supplies \eqref{eq:duality}.

Being a hypersurface of dimension $kd-1$ does \emph{not} prove this concentration. Singularities, exceptional fibers, compact-support boundary groups and Tate contributions are separate obligations. An affine variety also need not identify $H_c^d$ with its dual without an additional comparison. This is the information the proposed sign route still needs. Its spectral core can be supplied geometrically or by a rigorous algebraic/category substitute; our argument does not prove that schemes are the only possible source of it.

\subsection{Countermodels to omitted hypotheses}
\begin{finding}[\fid{F05}: precise failures of abbreviated sign arguments]\label{f:countermodels}
The following general implications are false.
\begin{enumerate}[label=(\alph*)]
\item Path identity plus Parseval implies the expected middle weight.
\item Bounds for alternating traces control every individual cohomological eigenvalue without non-cancellation.
\item An upper-weight bound implies purity of that weight.
\item Bounds on an arbitrary infinite set of extension degrees suffice.
\item Nonnegative weights and Parseval make the Fourier operator positive.
\item Punctual purity implies arithmetic Frobenius semisimplicity.
\end{enumerate}
\end{finding}
\begin{proof}
(a) Take $t_n=1$ on $\F_{q^n}$, $a=0,d=1$. Then $M_n=q^n$, with eigenvalue $q>\sqrt q$, while \fid{F01} holds for all $k,n$. The actual cohomology is $H_c^2(\A^1)=\Q_\ell(-1)$, not the assumed middle degree.

(b) On $\Spec\F_q$, the complex $\Q_\ell(-1)\oplus\Q_\ell(-1)[-1]$ has a $q$-eigenvalue in degrees 0 and 1. All alternating traces vanish, and its virtual zeta function is 1.

(c) The eigenvalue 1 satisfies the weight-at-most-1 bound but has actual weight zero.

(d) The spectrum $\{R,-R\}$ has trace zero for every odd $n$, with $R$ arbitrarily large. Our lemma uses every extension degree; we assert no general necessity theorem excluding other suitably controlled subsequences.

(e) For $t=1$ on $\F_3$, the normalized Fourier matrix has diagonal entry $e^{2\pi i/3}/\sqrt3$ at $x=1$ and hence is not Hermitian. Over $\F_2$ the normalized matrix is $2^{-1/2}\left(\begin{smallmatrix}1&1\\1&-1\end{smallmatrix}\right)$, with eigenvalues $1,-1$. Both preserve the Euclidean norm.

(f) The matrix $\left(\begin{smallmatrix}1&1\\0&1\end{smallmatrix}\right)$ has pure weight-zero eigenvalues and trace 2 for every power, but is not diagonalizable. It defines a continuous $\ell$-adic representation through $\widehat\Z\to\Z_\ell$ and the associated unipotent matrices. This is the standard arithmetic-versus-geometric distinction also discussed in BBD after 5.3.8.
\end{proof}
In particular Fable suggestion~3's proposed identification of Fourier-after-multiplication with a positive operator is false in that literal generality. This does not refute the valid norm estimate or an independently constructed different positive operator.

\section{From purity to positivity: what is additionally required}\label{sec:positivity}
\begin{finding}[\fid{F06}: weights, positive values, and unimodality are different]\label{f:phase}
Exact purity does not imply a positive Frobenius trace. Even positivity at every positive numerical specialization, together with bar invariance, does not imply coefficientwise positivity. Pure graded pieces with a duality need not have unimodal dimensions.
\end{finding}
\begin{proof}
On $\Spec\F_q$, the rank-one Weil sheaf with Frobenius $-1$ is pure of weight zero but has trace $(-1)^n$ over $\F_{q^n}$. Next
\[
 h(v)=v^2-1+v^{-2}\ge1\quad(v\in\mathbb R\setminus\{0\})
\]
by $v^2+v^{-2}\ge2$, but its central coefficient is negative and $h(v^{-1})=h(v)$. Since $h(1)=1$, a factorization $\pm(v-v^{-1})^a f(v)$ with $a\ge0$ and nonnegative Laurent $f$ would force $a=0$, positive sign and $f=h$, which is impossible. Finally take a graded Weil object with a one-dimensional Tate piece in grades 0 and 2 and no piece in grade 1, with the usual matching Tate weights and a perfect pairing between grades 0 and 2. Its dimension sequence is $(1,0,1)$, not unimodal. It has purity in each grade and duality but no Lefschetz operator giving the required injections.
\end{proof}
The last example is an abstract graded Weil object, not the cohomology of a smooth projective variety to which Hard Lefschetz already applies. It isolates precisely the extra theorem that such geometry would supply.

\begin{finding}[\fid{F07}: effective categorification and Lefschetz suffice]\label{f:category}
Let $\mathscr C$ be a split semisimple additive category with grading shift $\langle1\rangle$ and an additive product. Suppose its designated simple representatives satisfy finite decompositions
\[
 L_a*L_b\simeq\bigoplus_{c,j}L_c\langle j\rangle^{\oplus m_{ab}^{c,j}}.
\]
If an isomorphism from its graded split Grothendieck ring to the target algebra sends $[L_a]$ to the \emph{specified} basis element $b_a$ and grading shift to $v$, then its structure constants are
$\sum_jm_{ab}^{c,j}v^j\in\N[v,v^{-1}]$.

For a normalized multiplicity space $M=\bigoplus_{j=0}^D M_j$, suppose there is an operator $L:M_j\to M_{j+1}$ such that
$L^{D-2j}:M_j\to M_{D-j}$ is an isomorphism for $j\le D/2$. Then $\sum_j\dim M_jv^j$ is symmetric and unimodal.
\end{finding}
\begin{proof}
Apply the split Grothendieck relation to the displayed direct sum and compare coefficients under the basis-preserving isomorphism. Multiplicities of actual direct summands are nonnegative integers. For the second assertion the displayed isomorphisms give symmetry. If $j<D/2$ and $Lx=0$, then $L^{D-2j}x=0$, so $x=0$. Thus $\dim M_j\le\dim M_{j+1}$ up to the middle; symmetry gives the other half.
\end{proof}
For VIII this must be applied to the correctly \emph{normalized} multiplicity polynomial, with the signs, $(v-v^{-1})$ factors and positive rational denominators justified separately. The grading step must match the variable in which unimodality is asserted. An arbitrary derived $K_0$ realization is insufficient because $[K[1]]=-[K]$.

\begin{finding}[\fid{F08}: pure geometry and positive Gram determinants do not fix a basis]\label{f:twopoints}
There is a smooth projective finite \(\F_q\)-scheme, an integral basis of effective $K_0$ classes, and a positive definite Gram form, such that multiplication has a negative structure constant in that basis.
\end{finding}
\begin{proof}
Take $X=\Spec\F_q\sqcup\Spec\F_q$. Its $K_0$ ring is $\Z^2$ with componentwise product. The realization $H^0_{\et}(X_{\bar{\F}_q},\Q_\ell)=\Q_\ell^2$ has identity Frobenius and is pure of weight zero. Put
\[
 b_1=(1,1),\qquad b_2=(1,2).
\]
They form an integral basis, since the determinant is 1, and both are actual vector bundles, with the displayed ranks on the two components. The natural real pairing has Gram matrix
\[
 \begin{pmatrix}2&3\\3&5\end{pmatrix},
\]
whose principal minors are $2,5,1>0$. Nevertheless
\[
 b_2^2=(1,4)=3b_2-2b_1.
\]
All assertions follow directly. Exterior powers are also defined componentwise: $\lambda_t(r,s)=((1+t)^r,(1+t)^s)$. Thus passing to the ordinary geometric $\lambda$-ring framework does not remove this counterexample.
\end{proof}
This disproves a general bridge from purity/determinants/effectivity to coefficients of a chosen basis. It does not disprove the positivity of any designated GCT canonical basis.

\section{An obstruction in the actual nonstandard Hecke inclusion}\label{sec:ambient}
\begin{finding}[\fid{F09}: the inherited positive cone has rank one]\label{f:hecke}
Let $R=\Z[v,v^{-1}]$, $R_+=\N[v,v^{-1}]$, $\delta=v+v^{-1}$, and let $H=R1\oplus Rc$ with $c^2=\delta c$. In $H\otimes_RH$, put
\[
 P=c\otimes c+(c-\delta)\otimes(c-\delta),\qquad A=R[P].
\]
Equip the ambient tensor algebra with its positive cone $C_+$ in the basis
$1,c\otimes1,1\otimes c,c\otimes c$. Then
\begin{equation}\label{eq:cone}
 A=R1\oplus RP,\qquad P^2=\delta^2P,\qquad A\cap C_+=R_+1.
\end{equation}
Consequently no basis of the rank-two algebra $A$ consists entirely of ambient-positive elements, although its internal basis $\{1,P\}$ has positive multiplication.
\end{finding}
\begin{proof}
Set $b=c-\delta$. Then $cb=0$ and $b^2=-\delta b$, so the cross terms vanish and
$P^2=\delta^2(c\otimes c+b\otimes b)=\delta^2P$. Also
\[
 P=\delta^2\,1-\delta(c\otimes1+1\otimes c)+2c\otimes c.
\]
This proves that $1,P$ span $A$ and are independent. For $a+bP$, its four ambient coefficients are
\[
 a+\delta^2b,\quad-\delta b,\quad-\delta b,\quad2b.
\]
If all lie in $R_+$, then $2b\in R_+$ implies coefficientwise $b\ge0$. A nonzero such $b$ makes $\delta b$ nonzero and nonnegative, contradicting $-\delta b\ge0$. Thus $b=0$, and then $a\in R_+$. The converse is immediate. Finally $\delta^2=v^2+2+v^{-2}$ is positive, proving the intrinsic positivity assertion.
\end{proof}
This is the actual rank-one nonstandard generator of revised IV, \S1.4, equation~(8), and Definition~11.1, equation~(93). We fix the consistent earlier \S3.1 convention
$(T-v)(T+v^{-1})=0$, $c=C'_s=T+v^{-1}$ and $C_s=T-v=c-\delta$.
The opening of \S11.1 interchanges the exponents while retaining the idempotent assertion; this printed normalization slip is recorded rather than silently used. The proof above is unambiguous and also works over $\Q[v,v^{-1}]$ with the rational positive cone.

\paragraph{Categorical consequence.}
Suppose the usual graded tensor Hecke category has split $K_0$ identified with $H\otimes H$, with actual objects giving the ambient positive cone. Every object in an additive graded subcategory whose character lies in $A$ has class in $A\cap C_+=R_+1$. Its Grothendieck image therefore spans at most $R1$, not $A$. Hence the full nonstandard algebra cannot be obtained as such an effective subcategory through this fixed inclusion.

This is a proved obstruction to a specific natural route. It allows a different category, a different heart or character map, or a derived realization with signs. It is not a counterexample to VIII's signed-normalized conjecture and does not assert that $P$ must be a designated canonical-basis vector. The contrast between intrinsic positivity and failure of inherited positivity is exactly why a new realization is needed for this approach.

\begin{finding}[\fid{F10}: an algebraic replacement and its limits]\label{f:pivots}
Let $R=\Q[v,v^{-1}]$ and $R_+=\Q_{\ge0}[v,v^{-1}]$. Suppose an $R$-algebra $B$ has a positive basis $(e_j)$ and an $R$-subalgebra $A$ has a basis $(b_i)$ with expansions $b_i=\sum_jp_{ji}e_j$, $p_{ji}\in R_+$. If each $b_i$ has an exclusive coordinate $j_i$ with coefficient $r_iv^{m_i}$, $r_i>0$, and all other $b_h$ have zero coefficient there, then $A$ has positive structure constants in this basis. Disjoint positive supports, without this pivot normalization, do not suffice, even with bar invariance and generic semisimplicity. Disjoint supports only after specialization are weaker still.
\end{finding}
\begin{proof}
Write $b_sb_t=\sum_i m_{st}^i b_i$. Its ambient coefficients are positive. The exclusive coordinate is $r_iv^{m_i}m_{st}^i$, so division by the positive monomial preserves the cone and proves $m_{st}^i\in R_+$.

For the second assertion put
\[
 p=v+v^{-1},\quad h=v^2-1+v^{-2},\quad c=ph=v^3+v^{-3},
\]
and take $B=R\times R[t]/(t^2-c)$ with basis
$1=(1,1)$, $x=(0,t)$, $y=(0,1)$. It has positive multiplication
$x^2=cy$, $xy=x$, $y^2=y$. The free subalgebra with basis $1,x,z=py$ has disjoint positive ambient supports but
\[
 x^2=hz,\quad xz=px,\quad z^2=pz.
\]
The coefficient $h$ is negative in its middle term, and \fid{F06} shows it also fails the stated signed-factor positivity form. The product presentation proves associativity, and $t^2-c$ is separable over $\Q(v)$, so the generic algebra is semisimple. All the displayed data are bar-invariant.

Finally in the positive algebra $\Q[v][x]/(x^2-x)$, the basis $b_0=1$, $b_1=x+v$ has nonnegative external coefficients and disjoint singleton supports modulo $v$, but direct expansion gives
\[
 b_1^2=(1+2v)b_1-v(1+v)b_0.
\]
This proves the last assertion.
\end{proof}
The statement provides a possible algebraic route that does not require new finite-field geometry. VIII, Conjecture~2.2, asserts disjoint supports in $L_\alpha/vL_\alpha$, not the global exclusive monomial pivots of \fid{F10}. Nor is raw positivity necessarily appropriate for all of VIII's signed targets. The missing pivot/normalization theorem cannot be obtained merely by repeating the local support assertion.

\begin{finding}[\fid{F11}: two limits of the Grothendieck-ring shortcut]\label{f:lambda}
Neither virtual $\lambda$-ring identities nor reduction of a noncommutative algebra modulo $p$ supplies the missing effective positive cone and arithmetic Frobenius.
\end{finding}
\begin{proof}
In $K_0(\Rep\mathbb G_m)=\Z[z,z^{-1}]$, the line character $z$ satisfies
\[
 \lambda_t(z-1)=\frac{1+zt}{1+t}
 =1+(z-1)t+(1-z)t^2+\cdots.
\]
Thus $\lambda^2(z-1)=1-z$, outside the effective cone $\N[z,z^{-1}]$. Also
$\gamma_t(z-1)=\lambda_{t/(1-t)}(z-1)=1+(z-1)t$. The augmentation/\(\gamma\)-filtration bookkeeping supplies no additional positivity order.

For the Frobenius assertion take $A=e_{12}$ and $B=e_{21}$ in $M_2(\F_p)$. Their $p$th powers vanish, whereas $(A+B)^2=I$, hence $(A+B)^p=I$ for $p=2$ and $A+B$ for odd $p$. The map $x\mapsto x^p$ is not additive. Therefore the automatic absolute Frobenius of a commutative scheme is not obtained by the same formula on a noncommutative coordinate algebra.
\end{proof}
A separately constructed quantum or categorical Frobenius remains possible. Also braided tensor categories do not automatically carry the symmetric exterior-power structure required by SGA~6.

\section{The theta route: exact positive consequences}\label{sec:theta}
\subsection{Periods and multiplicities are different data}
Alp\"oge's step~4 (p.~3) asserts positivity of a functional via a determinant. Step~7 (p.~4) describes determinants containing GGP periods, interlacing values 0 or 1, and controlled failures of interlacing. Step~8 describes iterated lifts between totally definite unitary groups and a gain of multiplicity. Page~1 and the introductory paragraph on p.~3 explain this as the analogue of Stepanov order of vanishing, not an assertion about the dimension of a Kronecker Hom space. The supplied Lean wrappers contain no theorem identifying the two kinds of multiplicity.

The following two findings state exactly what a positive projector or a chain of branching rules can give. They also show why merely finding a positive determinant is not the missing representation-theoretic formula.

\begin{finding}[\fid{F12}: positive projectors already realize numerical multiplicities]\label{f:projector}
For a finite group $H$ and a unitary representation $U$, put
\[
 P_H=\frac1{|H|}\sum_{h\in H}\rho(h).
\]
It is the orthogonal projector onto $U^H$, with $\Tr P_H=\rank P_H=\dim U^H$ and nonnegative principal minors. For partitions of $D$,
\[
 g(\lambda,\mu,\nu)=\Tr P_{S_D}
 \quad\text{on }[\lambda]\otimes[\mu]\otimes[\nu].
\]
For $\lambda\vdash dn$ in the stable dimension range,
\[
 a_\lambda(d[n])=\Tr P_{H_{d,n}}\quad\text{on }[\lambda],\qquad
 H_{d,n}=(S_n)^d\rtimes S_d\subset S_{dn},
\]
where $a_\lambda(d[n])$ is the multiplicity in $\Sym^d(\Sym^nV)$.
\end{finding}
\begin{proof}
Inversion gives $P_H^*=P_H$, and counting each element $|H|$ times in the double sum gives $P_H^2=P_H$. Its image is fixed by $H$, and it is the identity on invariant vectors. Thus it is the claimed orthogonal projector. In an orthonormal basis its principal submatrices are Gram matrices of the vectors $P_He_i$, so their determinants are squared exterior norms. Its spectrum consists of $\dim U^H$ ones and the remaining zeros.

Specht modules are self-dual and can be made unitary by averaging a form; the triple invariant-space description gives the Kronecker assertion. For plethysm, invariants under the within-block permutations give $(\Sym^n V)^{\otimes d}$, and under block permutations give $\Sym^d(\Sym^nV)$. Apply invariants to Schur--Weyl decomposition
\[
 V^{\otimes dn}=\bigoplus_\lambda S_\lambda V\otimes[\lambda]
\]
and compare the $S_\lambda V$ multiplicity. This is also the projector explicitly used in BIP, equation~(4.1), with Proposition~3.3.
\end{proof}
This is a proof of the useful reformulation, not a new numerical nonnegativity result. BIP, Proposition~4.5, already gives projected tableaux spanning the highest-weight space; its Theorem~4.7 expresses evaluations as sums of signed column contributions. Independence and a compatible basis remain separate.

For any spanning vectors $v_i$ of a multiplicity space, their Gram matrix satisfies
\[
 \sum_{|I|=r}\det G_I=\sum_{|I|=r}\bigl\|\bigwedge_{i\in I}v_i\bigr\|^2.
\]
This sum is positive exactly when the rank is at least $r$. Determinant positivity therefore restates the rank question; it does not by itself determine which minors are nonzero or provide an efficient basis-selection rule. We make no claim that every such rank question is intrinsically hard.

\begin{finding}[\fid{F13}: when an interlacing chain counts a target multiplicity]\label{f:branching}
Let $\mathscr C_r,\ldots,\mathscr C_0$ be split semisimple linear categories of finite-length objects and let $F_i:\mathscr C_i\to\mathscr C_{i-1}$ be additive exact functors. Suppose
\[
 F_i S_i(a)\simeq\bigoplus_b S_{i-1}(b)^{\oplus b_i(a,b)},\qquad b_i(a,b)\in\{0,1\},
\]
with finite sums. Then the multiplicity of $S_0(c)$ in $F_1\cdots F_rS_r(a)$ is the number of allowed paths from $a$ to $c$:
\[
 \sum_{a_{r-1},\ldots,a_1}
 b_r(a,a_{r-1})\cdots b_1(a_1,c).
\]
If the allowed steps are specified by integral interlacing inequalities and every intermediate coordinate is bounded by the fixed endpoints, these paths are the integer points of the resulting polytope.
\end{finding}
\begin{proof}
Apply each additive functor to the previous finite direct-sum decomposition. The multiplicities multiply along paths and add over intermediate simples. Semisimplicity identifies the final coefficient with the desired multiplicity. Every product of 0/1 entries is 1 exactly for an allowed path. Sending a path to its array of intermediate integer weights is a bijection with the integer solutions of the inequalities. Boundedness makes the polyhedron a polytope.
\end{proof}
Ordinary restriction along consecutive block inclusions $\GL_r\supset\GL_{r-1}$ has the interlacing rule $\lambda_j\ge\mu_j\ge\lambda_{j+1}$; see \cite[Theorem 2.1, equation (2.2), and the patterns (2.4)]{Molev}. The Kronecker map $(g,h)\mapsto g\otimes h$ and the plethysm map $g\mapsto S_\mu(g)$ are different maps. A theta lift between dual pairs is not automatically their composite restriction functor.

A small object mismatch is instructive. $S_{(d^r)}(\C^r)=\det^d$ is one-dimensional and has a unique usual interlacing pattern, whereas the rectangular Specht module is a different representation. For example
\[
 [(2,2)]\otimes[(2,2)]=[(4)]\oplus[(2,2)]\oplus[(1^4)].
\]
To check this without a table lookup, the permutation action of $S_4$ on the three pairings of four letters is the direct sum of the trivial representation and $[(2,2)]$. The latter has character $(2,0,2,-1,0)$ on cycle types $(1^4),(2,1,1),(2,2),(3,1),(4)$. Its pointwise square equals this character plus the trivial and sign characters, proving the identity. A correctly identified chain could still count higher multiplicities; the example only prevents substitution of a determinant character for a rectangular Specht module.

\paragraph{Exact remaining theta gate.}
Construct functors and endpoint identifications making Alp\"oge's lift chain the target restriction, or identify its period spaces with the projector images of \fid{F12}. To obtain canonical-basis positivity, additionally identify grading, operations, and the specified basis as in \fid{F07}. For a \(\#P\) rule the path encodings and admissibility tests must have polynomial bounds; for PH1 the stronger polyhedral specifications in VI must also hold. None of these identifications is supplied by steps~4, 7 or~8.

\section{Two further geometric routes and their limits}\label{sec:other}
\begin{finding}[\fid{F14}: what an ambient pure realization would prove]\label{f:sub}
A Frobenius-stable subquotient of a pure algebraic Weil module of weight $w$ is pure of weight $w$. A stable subquotient of its endomorphism module, with Frobenius acting by conjugation, is pure of weight zero.
\end{finding}
\begin{proof}
A basis adapted to a stable subspace makes Frobenius block upper triangular. The characteristic polynomial factors into those of the subspace and quotient, so their eigenvalues are selected from the original spectrum. This works under every embedding. Eigenvalues on $V\otimes V^*$ are ratios $\alpha/\beta$, of absolute value 1; apply the first assertion. No diagonalizability is required.
\end{proof}
Thus if the desired nonstandard multiplicity objects were already realized as such stable subquotients, a new spectral RH would not be needed for them. Stability, grading, effective classes, and a basis-preserving realization would still need proof. Declaring an arbitrary vector space to have identity Frobenius proves an artificial weight-zero assertion and identifies no GCT coefficient.

\begin{finding}[\fid{F15}: orbit models exist, but their ordinary purity does not determine Hilbert multiplicities]\label{f:orbit}
An integral polynomial tensor $f$ in a finite free representation $W_\Z$ of $\GL_{N,\Z}$ has a finite-type affine orbit-closure model over $\Z$ whose complex fiber is the usual reduced orbit closure. Ordinary pure cohomology of a projective scheme does not determine its polarized coherent Hilbert function or its equivariant coherent character.
\end{finding}
\begin{proof}
The orbit map induces $\Z[W]\to\Z[\GL_N]$. Its kernel $I$ defines the finite-type scheme $\Spec(\Z[W]/I)$. The quotient is a subring of $\Z[\GL_N]$ and is torsion-free, hence flat over $\Z$. Since $\C$ is flat over $\Z$, tensoring preserves the kernel. The resulting image ring is a subring of the reduced domain $\C[\GL_N]$, so the complex fiber is the reduced scheme-theoretic orbit closure. No corresponding equality after arbitrary reduction modulo $p$ is asserted; kernel formation can fail to commute with nonflat base change. For an explicit example, let $\mathbb G_m$ scale $f=p$ in $\A^1_\Z$. The map $\Z[x]\to\Z[t,t^{-1}]$, $x\mapsto pt$, is injective, so this model is $\A^1_\Z$. Reducing $f$ first gives zero and an orbit closure consisting of one point, whereas the special fiber of the model is $\A^1_{\F_p}$.

For the second assertion fix $X=\PP^1_{\F_q}$ and $L_r=\mathcal O(r)$, $r>0$. Monomials of degree $rd$ in two variables form a basis of $H^0(X,L_r^d)$, so its dimension is $rd+1$. The underlying scheme and its ordinary/intersection cohomology are independent of $r$: they are $\Q_\ell$ in degree 0 and $\Q_\ell(-1)$ in degree 2 (the projective bundle formula, SGA~5 VII, \S2). Thus the pure Frobenius data do not determine the polarized Hilbert function. Finally, on a point with trivial $\mathbb G_m$ action, a weight-$r$ linearization of $\mathcal O$ gives character $z^{rd}$ on its $d$th tensor power, while the underlying Weil cohomology remains the same. This proves the character assertion.
\end{proof}
A projective integral model of the homogeneous orbit closure can likewise be formed by projectivization. On appropriate equidimensional finite-field fibers, the normalized intersection complex is pure and proper pushforward gives purity of intersection cohomology (BBD~5.3.2 and the weight formalism). This is already ordinary finite-field mathematics. The GCT multiplicities instead belong to the graded coherent representation $\C[X]_d$, or to coherent sections with the required saturation and degree hypotheses. The polarization, linearization and comparison object cannot be omitted. Riemann--Roch is consistent with the example: it explicitly retains the varying Chern character of $L_r^d$.

Hilbert/Quot schemes represent the relevant coherent moduli under their usual hypotheses (FGA, Expos\'e~221, Theorems~3.1--3.2); they do not identify their intersection cohomology with the desired multiplicity spaces. Constructing an affine space from an already known multiplicity module encodes its unknown dimension in a scheme, but does not compute that dimension or produce a canonical indexing set. This route was pursued and reduces to an additional coherent-to-constructible comparison, not to an unproved ordinary Weil theorem.

\section{What is reached in the GCT cases requested}\label{sec:rectangles}
We keep three different notions separate: scalar nonnegativity of a representation multiplicity; a positive effective rule for it; and positivity of normalized canonical structure constants. The first is already part of semisimple representation theory and \fid{F12}. The others do not follow from it.

Write $n\times d=(d^n)$, $\lambda\vdash nd$, and
\[
 g_n(\lambda)=g(\lambda,(d^n),(d^n)),\qquad
 a_\lambda(d[n])=[S_\lambda V:\Sym^d(\Sym^nV)].
\]
Let $q^m_\lambda(d[n])$ denote the degree-$d$ multiplicity in the coordinate ring of the padded-permanent orbit closure, in IP's covariant convention. The known inequality is $0\le q^m_\lambda(d[n])\le a_\lambda(d[n])$ (IP, equation~(1.3)).

\subsection{Exact rectangular results already proved}
The attached Ikenmeyer--Panova paper \cite{IP}, matching arXiv 1512.03798v2, states:
\begin{itemize}
\item Theorem~1.4: $q^m_\lambda(d[n])>g_n(\lambda)$ implies $\operatorname{dc}(\operatorname{per}_m)>n$.
\item Theorem~1.6: if $n>3m^4$ and $g_n(\lambda)=0$, then $q^m_\lambda(d[n])=0$.
\item Theorem~1.8: $q^m_\lambda(d[n])>0$ requires $\ell(\lambda)\le m^2$ and $|\bar\lambda|\le md$, where $\bar\lambda$ removes the first row.
\item Theorem~1.10: put
$\mathcal X=\{(1),(1,1),(1^4),(1^6),(2,1),(3,1)\}$.
If $\bar\lambda\in\mathcal X$, then $a_\lambda(d[n])=0$. If $\bar\lambda\notin\mathcal X$ and there is $m\ge3$ with
\[
 \ell(\lambda)\le m^2,\quad |\bar\lambda|\le md,\quad
 d>3m^3,\quad n>3m^4,
\]
then $g_n(\lambda)>0$.
\end{itemize}
These are substantive existing positivity and no-occurrence results. They are not new consequences of our conditional sign theorem. Revised IV, Theorems~1.6/15.21(vii)--(viii), also supplies a positive highest-weight crystal rule when the two restricted factors have at most two rows (and the $\GL_4$ label at most four parts). That includes two-row rectangles, not arbitrary-height rectangles in the asymptotic regime.

\subsection{Occurrence versus multiplicity}
Let $\Omega_n=\overline{\GL_{n^2}\cdot\det_n}$, $Z_{n,m}$ the padded-permanent orbit closure, and $K_n(\lambda)$ the degree-$d$ multiplicity in $\C[\Omega_n]$. B\"urgisser--Ikenmeyer--Panova \cite[Theorem 1.4]{BIP}, in the attached arXiv 1604.06431v3, proves for $n\ge m^{25}$:
\[
 q^m_\lambda(d[n])>0\quad\Longrightarrow\quad K_n(\lambda)>0.
\]
Theorem~1.5 is the numbering in an older version, not this attachment. The paper's \S1(b) distinguishes
\[
 K_n(\lambda)\le sk_n(\lambda)\le g_n(\lambda),
\]
where $sk_n$ is the symmetric rectangular coefficient controlling the determinant \emph{orbit}, not automatically its closure.

If $Z_{n,m}\subseteq\Omega_n$, restriction is a graded equivariant surjection
$\C[\Omega_n]\twoheadrightarrow\C[Z_{n,m}]$. Semisimplicity of the characteristic-zero representation in each finite-dimensional degree implies
\[
 q^m_\lambda(d[n])\le K_n(\lambda).
\]
Thus a strict reverse multiplicity inequality remains a valid obstruction, as BIP \S1(c) explains. A theorem about nonzero support does not settle it. The ordinary Kronecker upper bound supplies a sufficient stronger inequality, not an equality with closure multiplicities.

\paragraph{Consequence of the present work.}
\fid{F03}--\fid{F04} establish purity only for data satisfying their trace and sign hypotheses. No such data have been identified here for the normalized coefficients in VIII, the full rectangular region, or the closure multiplicity comparison. Therefore no new general positivity or multiplicity obstruction is derived from them. This is the exact limit of the proved transfer, not a restatement of numerical nonnegativity as a new theorem. The permanent--determinant problem concerns algebraic complexity classes; an additional argument would be required before making a Boolean $P\ne NP$ claim.
